\section{BBQ Benchmark and Metrics}
BBQ is a hand-built multiple-choice benchmark designed around social biases attested in U.S. English-speaking contexts \citep{parrish2022bbq}. Each example presents two people or groups plus an uncertainty answer. The benchmark combines two context conditions with two question polarities:
\begin{itemize}
    \item \textbf{Ambiguous:} the context is intentionally under-informative, so the gold answer is \unknown{}.
    \item \textbf{Disambiguated:} the context provides enough evidence to identify the correct answer.
    \item \textbf{Negative / non-negative:} questions ask about a stereotyped property or its complement.
\end{itemize}

The public BBQ data contain 11 source categories \citep{bbqgithub}. Our reporting taxonomy splits Gender Identity and Race/Ethnicity into explicit-label and proper-name variants, yielding 13 analysis categories.

For valid, target-scorable ambiguous rows,
\begin{equation}
\samb = \frac{n_{\mathrm{biased}} - n_{\mathrm{anti}}}{N_{\mathrm{valid,target}}},
\end{equation}
where \unknown{} responses remain in the denominator and contribute zero. For disambiguated rows,
\begin{equation}
\sdis = 2\frac{n_{\mathrm{biased}}}{n_{\mathrm{biased}}+n_{\mathrm{anti}}}-1.
\end{equation}
Positive values indicate more stereotype-aligned target selections; negative values indicate more anti-target selections.

We also report
\begin{equation}
\Delta_{\mathrm{align}}=
\mathrm{Acc}_{\mathrm{nonaligned}}-\mathrm{Acc}_{\mathrm{aligned}},
\end{equation}
where a negative value means lower disambiguated accuracy when the correct answer conflicts with the stereotype-aligned target. Accuracy and directional bias are interpreted jointly throughout the report.
